性厌恶（sexual aversion）指对性接触产生强烈的厌恶、回避与焦虑反应，甚至一想到或一接触性相关情境（拥抱、接吻、性交）就出现恶心、心悸、出汗、发抖等躯体反应。它与"暂时不想做"的性欲低下不同，是一种带有明显回避与负性情绪的性心理问题。

\begin{itemize}
	\item \textbf{常见成因}：性创伤或性侵经历、童年期不良经历、严苛或羞耻化的性教育、被强迫或被贬低的性经历、对体液的过度厌恶、伴侣关系中的控制或暴力、焦虑与强迫特质。
	\item \textbf{与阴道痉挛的关系}：性厌恶常与阴道痉挛、性交疼痛恐惧并存，互为因果——恐惧引发肌肉紧张与疼痛，疼痛又强化厌恶。
	\item \textbf{并非"不爱对方"}：性厌恶是针对"性情境"的强烈反应，而非对伴侣情感的否定；把它解读为不爱，只会加重双方的痛苦。
	\item \textbf{应对路径}：停止强迫与"硬闯"；在安全、自愿、可控的前提下，从非性的亲密（牵手、拥抱、按摩、依偎）逐步重建身体安全感；配合认知行为治疗、系统脱敏与盆底放松训练；既往有创伤者需采用创伤知情（trauma-informed）的心理治疗。
	\item \textbf{伴侣的角色}：理解、耐心与"随时可以停止"的边界，是康复的前提；勉强与施压只会加深厌恶。
	\item \textbf{何时就医}：若厌恶强烈、持续、已明显影响关系与情绪，或与创伤经历有关，建议寻求性治疗师或心理科的专业帮助——这是可以治疗的。
\end{itemize}

\noindent\textit{【编者注】}性厌恶不是"矫情"或"不够爱"，而是真实的性心理问题，且对治疗有反应。若伴随闪回、噩梦、回避等创伤后应激表现，应优先进行专业的创伤评估与治疗。
